\documentclass[10pt]{article}
\usepackage[T1]{fontenc}
\usepackage[utf8]{inputenc}
\usepackage[french]{babel}
\usepackage{amsmath,amssymb,mathrsfs,mathpazo}
\usepackage{geometry}
\geometry{a4paper,margin=2cm}
\usepackage{graphicx,array,tabularx,tabularray,colortbl,multirow,multicol,booktabs,enumitem,pifont,ulem,fancybox,diagbox,cancel}
\usepackage{pstricks,pstricks-add,pst-plot,pst-tree,pst-node,pst-text,pst-func,pst-math,pst-bspline,pst-3dplot,pst-eucl,pst-solides3d}
\let\pscalculate\undefined
\usepackage{pst-bezier}
\usepackage{tikz,pgfplots,tkz-tab}
\usetikzlibrary{arrows,arrows.meta,patterns,calc,positioning,trees,decorations.markings,shapes.geometric,shapes.misc,shapes.arrows,intersections,decorations.pathmorphing,decorations.pathreplacing,backgrounds,shadows,babel}
\pgfplotsset{compat=1.18}
\usepgfplotslibrary{fillbetween}
\usepackage[np]{numprint}
\usepackage{listings,makecell,pas-tableur,twoopt}
\let\barre\undefined
\usepackage[pstricks]{tablvar}
\definecolor{darkgreen}{RGB}{0,100,0}
\providecommand{\red}{\color{red}}
\providecommand{\blue}{\color{blue}}
\providecommand{\green}{\color{green}}
\providecommand{\black}{\color{black}}
\providecommand{\white}{\color{white}}
\providecommand{\gray}{\color{gray}}
\providecommand{\cyan}{\color{cyan}}
\providecommand{\magenta}{\color{magenta}}
\providecommand{\yellow}{\color{yellow}}
\providecommand{\violet}{\color{violet}}
\providecommand{\orange}{\color{orange}}
\providecommand{\euro}{\texteuro}
\providecommand{\EUR}{\texteuro}
\providecommand{\ieme}{\textsuperscript{e}}
\providecommand{\degres}{\ensuremath{^\circ}}
\providecommand{\degree}{\ensuremath{^\circ}}
\providecommand{\vv}[1]{\overrightarrow{#1}}
\providecommand{\Coord}[1]{#1}
\providecommand{\decosix}{}
\providecommand{\starredbullet}{\ensuremath{\bullet}}
\pagestyle{empty}
\setlength{\parindent}{0pt}
\usepackage{adjustbox}\DeclareUnicodeCharacter{2236}{:}\DeclareUnicodeCharacter{2190}{\ensuremath{\leftarrow}}
\providecommand{\R}{\mathbb{R}}
\providecommand{\N}{\mathbb{N}}
\providecommand{\Z}{\mathbb{Z}}
\providecommand{\Q}{\mathbb{Q}}
\providecommand{\D}{\mathbb{D}}
\providecommand{\C}{\mathbb{C}}
\providecommand{\e}{{\mathrm{e}}}
\providecommand{\vect}[1]{\overrightarrow{#1}}
\providecommand{\vectt}[1]{\overrightarrow{\mathrm{#1}}}
\providecommand{\vvt}[1]{\overrightarrow{\mathrm{#1}}}
\providecommand{\np}[2][]{{#2~#1}}
\providecommand{\vv}[1]{{\overrightarrow{#1}}}
\renewcommand{\d}{{\mathrm{d}}}
\providecommand{\degres}{{^\circ}}
\providecommand{\degre}{{^\circ}}
\providecommand{\blue}{\color{blue}}
\providecommand{\red}{\color{red}}
\providecommand{\enskip}{\quad}
\providecommand{\llbracket}{\lbrack\!\lbrack}
\providecommand{\rrbracket}{\rbrack\!\rbrack}
\providecommand{\euro}{{\text{€}}}
\providecommand{\no}{{\text{n°}}}
\providecommand{\interval}[2]{{[#1;#2]}}
\providecommand{\intervFO}[2]{{[#1;#2[}}
\providecommand{\dotfill}{\ldots}

\providecommand{\ds}{\displaystyle}
\providecommand{\textding}[1]{#1}
\providecommand{\Oij}{\(\mathrm{O};\vec{\imath},\vec{\jmath}\)}
\providecommand{\Oijk}{\(\mathrm{O};\vec{\imath},\vec{\jmath},\vec{k}\)}
\providecommand{\Ouv}{\(\mathrm{O};\vec{u},\vec{v}\)}
\providecommand{\intervFF}[2]{\ensuremath{[#1;#2]}}
\renewcommand{\euro}{\ifmmode\text{\texteuro}\else\texteuro\fi}
\DeclareUnicodeCharacter{00B0}{\ifmmode^{\circ}\else\textdegree\fi}
\DeclareUnicodeCharacter{20AC}{\ifmmode\text{\texteuro}\else\texteuro\fi}
% Énoncés : Éducation nationale. Transcription et corrigés : APMEP, Denis Vergès et contributeurs. Découpage et classement : CoopMaths, Sébastien Lozano, Éric Elter et contributeurs. Adaptation de la présentation : Maths Explorer.
% CC BY-SA : https://coopmaths.fr/www/about/
% https://coopmaths.fr/alea/static/e3c/2020/tex/e3c_2020_00_sujet9_4_cor.tex
\begin{document}
\noindent\textit{APMEP / CoopMaths. CC BY-SA. Présentation adaptée par Maths Explorer.}\par\medskip
\section*{Exercice 4 (5 points)}
\subsection*{1. On note $f'$ la fonction dérivée de $f$.}
\subsubsection*{a. Montrer que, pour tout réel $x$, $f'(x) = 3\left(x + \dfrac{1}{3}\right)(x - 1)$.}
	
	La fonction polynôme $f$ est dérivable sur $\mathbb{R}$ et donc sur $[0; +\infty[$ :
\begin{align*}
	f'(x)& = 3x^2 - 2x - 1 \\
	&= 3(x^2 - \dfrac{2}{3}x) - 1 \\
	&= 3\left[\left(x - \dfrac{1}{3}\right)^2 - \dfrac{1}{9}\right] - 1\\
&= 3\left[\left(x - \dfrac{1}{3}\right)^2 - \dfrac{4}{9}\right]\\
& = 3\left(x + \dfrac{1}{3}\right)(x - 1)
\end{align*}
On pouvait aussi développer l'expression proposée dans l'énoncé et prouver qu'elle est bien égale à la dérivée obtenue sous forme développée.
	\subsubsection*{b. En déduire le tableau de variation de $f$ sur $[0; +\infty[$.}
	
	Comme pour $x > 0$, $3\left(x + \dfrac{1}{3}\right) > 0$, le signe de $f'(x)$ est celui de $x - 1$, donc :
	\begin{itemize}
		\item Si $0 \leq x < 1$, $f'(x) < 0$ : la fonction est décroissante sur $[0; 1[$ ;
		\item Si $x > 1$, $f'(x) > 0$ : la fonction est croissante sur $[1; +\infty[$ ;
		\item $f(1) = 1 - 1 - 1 - 1 = -2$ est le minimum de la fonction sur $[0; +\infty[$.
	\end{itemize}
	
	\begin{center}
			\begin{tikzpicture}[double distance=2pt]
			\tkzTabInit{$x$/1,$f'(x)$/1,$f(x)$/2}{$0$,$1$, $+\infty$}
			\tkzTabLine{,-,z, +}
			\tkzTabVar{+/,-/ $-2$ /, +/}
		\end{tikzpicture}
	\end{center}
	\subsubsection*{c. Déterminer l’abscisse du point de la courbe représentative de $f$ pour lequel le coefficient directeur de la tangente vaut 7.}
	
Il faut trouver le nombre réel $x$ de $[0; +\infty[$ tel que :
	\[
	f'(x) = 7 \iff 3x^2 - 2x - 1 = 7 \iff 3x^2 - 2x - 8 = 0
	\]
	
	\[
	\Delta = 4 + 4 \times 3 \times 8 = 4 + 96 = 100 = 10^2 > 0
	\]
Cette équation a deux solutions :
	\[
	x_1 = \dfrac{2 + 10}{6} = 2 \quad \text{et} \quad x_2 = \dfrac{2 - 10}{6} = -\dfrac{4}{3}
	\]
Seule la première solution est positive, donc $f'(2) = 7$.
	
	\subsection*{2. On note $x_0$ l’unique solution de l’équation $f(x) = 0$. On admet que $x_0 \in [1; 2]$.}
	
	\subsubsection*{a. On considère la fonction suivante définie en langage Python.}
	
	\begin{center}
\begin{tabularx}{\linewidth}{|*{6}{>{\centering \arraybackslash}X|}}\hline
	Itération	&$x = \dfrac{a+b}{2}$	&$f(x)< 0$ ?&$a$&$b$&Amplitude de $[a~;~b]$\\ \hline
	$k = 0$		&1,5					&OUI		&1,5& 2	&0,5 \\ \hline
	$k = 1$		&1,75						&	OUI		&1,75	&2	&0,25\\ \hline
	$k = 2$		&	1,875					&	NON		&1,75	&1,875	&0,125\\ \hline
\end{tabularx}
\end{center}
	
	\subsubsection*{b. En déduire un encadrement de $x_0$, d’amplitude 0,125, par deux nombres décimaux.}
	

	
	On a donc $1,75 < x_0 < 1,875$.

\end{document}
